\documentclass[twocolumn]{aastex702}
\usepackage{amsmath}
\usepackage{booktabs}
\shorttitle{Common-denominator optical tracer census in SDSS}
\shortauthors{NebulaMind local integration}
\begin{document}

\title{Common-denominator optical tracer census in SDSS: selection-aware SDSS optical proxy integration}
\author{NebulaMind Research Autopilot}
\email{autopilot@nebulamind.com}
\affiliation{Local reproducible integration run; public SDSS DR17 data only}

\begin{abstract}
We integrate the active proposal 'A multiphase, common-denominator census of AGN-driven outflows' as a guarded SDSS DR17 optical denominator/proxy draft rather than as a completed physical-feedback paper. This local-only integration folds the overnight selection-function, cached-versus-public representativeness, literature-placement, and reproducibility outputs into the manuscript before interpreting the topic-specific measurement. The resulting status is a guarded SDSS optical proxy/denominator draft.
\end{abstract}

\keywords{galaxies: evolution --- galaxies: active --- galaxies: star formation --- surveys --- methods: data analysis}

\section{Purpose and claim contract}\label{sec:purpose}
This draft preserves the active proposal title, 'A multiphase, common-denominator census of AGN-driven outflows', but narrows the supported claim to the cached SDSS optical measurement named in the results. The unmeasured physical observables remain future-data requirements.

The claim contract is intentionally conservative. Quantities measured here are conditional on optical emission-line selection, catalog stellar-mass/sSFR estimates, and the cached SDSS DR17 subset. Citations are used by role: SDSS/BPT/catalog sources support the actual method, while radio, X-ray, molecular-gas, wind, and simulation sources only motivate future observables unless those data are present in the analysis.


\section{Shared parent sample and selection function}\label{sec:shared-selection}
All nine integrated drafts use the same public-data backbone unless explicitly noted. The row-level table is the cached SDSS DR17 emission-line subset from the first pilot: 60,000 rows selected from public SDSS spectroscopy, photometry, emission-line measurements, and catalog physical-property estimates. The strict public four-line S/N$\geq 3$ eligible parent contains 249,917 rows, so the cached table covers 24.0\% of that strict parent. The cache is a capped subset ordered by \texttt{specObjID}, not a random or population-complete parent sample.

\begin{deluxetable*}{lrrr}
\tabletypesize{\scriptsize}
\tablecaption{Shared SDSS DR17 selection cascade used before paper-specific quantities.\label{tab:selection-cascade}}
\tablehead{\colhead{Selection stage} & \colhead{Public DR17 rows} & \colhead{Cached rows} & \colhead{Retention vs. spectro-z parent}}
\startdata
SpecObj GALAXY, 0.02<z<0.12 & 501,060 & -- & 1.000 \\
plus galSpecInfo/PhotoObj/galSpecExtra and mass/sSFR bounds & 416,554 & -- & 0.831 \\
plus galSpecLine join & 416,554 & -- & 0.831 \\
four BPT lines positive with positive errors & 373,445 & 60,000 & 0.745 \\
four BPT lines S/N$\geq$3 & 249,917 & 60,000 & 0.499 \\
four BPT lines S/N$\geq$5 & 176,523 & 42,446 & 0.352 \\
four BPT lines S/N$\geq$10 & 91,768 & 22,311 & 0.183 \\
\enddata
\tablecomments{Counts are read-only public SDSS DR17 count queries plus the cached local CSV. Cached rows are shown only where the cache applies.}
\end{deluxetable*}

The four-line requirement is strongly selection dependent. In the public counts, S/N$\geq3$ keeps 33.6\% of the $-12<\log {\rm sSFR}<-11$ parent bin but 94.9\% of the $-10<\log {\rm sSFR}<-9.5$ bin. Therefore every incidence, quenching, density, gas-denominator, or target-vector statement in this integration is conditional on the four-line emission-line selection.

Cached-versus-public marginal checks found no redshift, stellar-mass, or sSFR bin with a cached-minus-public fraction difference above 5 percentage points. The largest absolute differences were 2.03 percentage points in redshift, -1.63 percentage points in stellar mass, and -0.58 percentage points in sSFR. This is a representativeness diagnostic only; it does not make the capped cache random or complete.


\section{Measurements}\label{sec:measurements}
The row-level measurements include redshift, stellar mass, catalog specific star-formation rate, model colors, and four optical line fluxes/errors. BPT labels and all derived denominators are recomputed locally from the cached table. Catalog SFR/sSFR values are treated as low-redshift SDSS physical-property estimates rather than direct resolved gas or feedback measurements \citep{sdssdr17,brinchmann2004,york2000}.


\section{Topic-specific optical denominator or proxy result}\label{sec:topic-result}
The consolidated proposal question is: How strongly do simple optical tracer definitions change the inferred AGN/feedback-candidate prevalence in one common SDSS denominator? The integrated manuscript deliberately demotes the claim to the directly measured SDSS quantity. It is not a full physical-feedback test.

\begin{itemize}
\item Within the same 60,000-galaxy denominator, simple optical tracer definitions produce prevalence from 0.136 to 0.418.
\item The widest-to-narrowest prevalence ratio is 3.1, before adding molecular, neutral, or X-ray/radio phases.
\item This demonstrates why a common-denominator multiphase census is required; it does not measure molecular or neutral outflow rates.
\end{itemize}


\begin{figure}
\centering
\includegraphics[width=\columnwidth]{../figures/fig-topic.pdf}
\caption{Topic-specific SDSS DR17 optical denominator/proxy diagnostic for debate-map-to-wiki-rebuild p1. The caption is intentionally narrow: this figure summarizes the cached optical result used for target definition or denominator design, not the unmeasured multi-survey physical claim.}
\label{fig:topic}
\end{figure}

\section{Interpretation and missing observables}\label{sec:missing}
SDSS-only pilot; full proposal requires the additional survey data named in the research-topic page. The full proposal requires: ionized, molecular, neutral, and X-ray/radio tracers measured over the same parent denominator and aperture model.

A real multiphase census needs independent ionized, neutral, molecular, and energetic outflow observables; the present SDSS thresholds are an optical denominator only \citep{veilleux2005,rupke2018,cicone2014,fiore2017,feruglio2015,woo2016,bae2018}.


\subsection{Literature Context and Missing Observables}

An optical ionized-gas census is tracer- and definition-dependent and cannot recover neutral/molecular mass loading \citep{veilleux2005,cicone2014,rupke2018}. Radio detection can change observed ionized-outflow incidence, motivating radio follow-up without changing this optical prevalence \citep{escott2025}. JWST observations at cosmic noon found neutral-gas outflow rates that can match or exceed ionized-gas rates, motivating a separate neutral tracer in any future common-denominator census \citep{davies2024}. Beam smearing and unresolved rotation block galaxy-wide outflow claims from a single aperture \citep{holden2024}.

\section{Reproducibility and safety}\label{sec:repro}
This manuscript was generated by local integration run \texttt{INTEGRATED\_9\_PAPERS\_20260709T012051Z}. Inputs are the original RP-1 SDSS query/run directory, the eight-topic SDSS remaining-topic manifest, the overnight shared selection-function packet, the cached-versus-public representativeness packet, Goru robustness outputs, and literature/source placement packets. The output is a local draft PDF and manifest entry only. No public-linked PDF was replaced.

\section{Conclusion}\label{sec:conclusion}
The integration improves the paper package by putting denominator honesty before results. For RP-1, the strongest outcome is a plausible short-paper association draft: broad optical BPT AGN hosts in this capped SDSS emission-line subset have lower catalog sSFR than mass--redshift matched star-forming controls, with robustness caveats. For the other active topics, the correct packaging is a guarded denominator/proxy suite, not eight independent causal feedback papers.


\begin{thebibliography}{99}
\bibitem[Abdurro'uf et al.(2022)]{sdssdr17} Abdurro'uf, Accetta, K., Aerts, C., et al. 2022, ApJS, 259, 35
\bibitem[Brinchmann et al.(2004)]{brinchmann2004} Brinchmann, J., Charlot, S., White, S.~D.~M., et al. 2004, MNRAS, 351, 1151
\bibitem[York et al.(2000)]{york2000} York, D.~G., Adelman, J., Anderson, J.~E., Jr., et al. 2000, AJ, 120, 1579

\bibitem[Bae \& Woo(2018)]{bae2018} Bae, H.-J., \& Woo, J.-H. 2018, ApJ, 853, 185
\bibitem[Cicone et al.(2014)]{cicone2014} Cicone, C., Maiolino, R., Sturm, E., et al. 2014, A\&A, 562, A21
\bibitem[Feruglio et al.(2015)]{feruglio2015} Feruglio, C., Fiore, F., Carniani, S., et al. 2015, A\&A, 583, A99
\bibitem[Fiore et al.(2017)]{fiore2017} Fiore, F., Feruglio, C., Shankar, F., et al. 2017, A\&A, 601, A143
\bibitem[Rupke(2018)]{rupke2018} Rupke, D.~S.~N. 2018, Galaxies, 6, 138
\bibitem[Veilleux et al.(2005)]{veilleux2005} Veilleux, S., Cecil, G., \& Bland-Hawthorn, J. 2005, ARA\&A, 43, 769
\bibitem[Woo et al.(2016)]{woo2016} Woo, J.-H., Bae, H.-J., Son, D., \& Karouzos, M. 2016, ApJ, 817, 108
\bibitem[Escott et al.(2025)]{escott2025} Escott, E.~L., Morabito, L.~K., Scholtz, J., et al. 2025, MNRAS, 536, 1166
\bibitem[Davies et al.(2024)]{davies2024} Davies, R.~L., Belli, S., Park, M., et al. 2024, MNRAS, 528, 4976
\bibitem[Holden \& Tadhunter(2024)]{holden2024} Holden, L.~R., \& Tadhunter, C.~N. 2024, MNRAS, 536, 1857
\end{thebibliography}

\end{document}
